% ============================================================
% KEY EQUATIONS FOR IEEE TRANS. SMART GRID (Journal Extension)
% Energy-Information Co-Optimization via Learnable Coupling
% ============================================================
% VERIFIED AGAINST SOURCE CODE: src/models/*.py
% Updated: 2026-02 — Journal version with auto-scaled K, 5 IEEE cases

% ============================================================
% SECTION: THEORETICAL FOUNDATION
% ============================================================

% ------------------------------------------------------------
% THEOREM 1: Delay-Stability Coupling Bound
% ------------------------------------------------------------
% WHAT: The central theoretical contribution establishing a formal
%       relationship between communication delay and power system stability.
% WHY:  Enables prediction of stability degradation from network latency,
%       allowing proactive grid management without trial-and-error.
% WHERE: Applied in the loss function to enforce stability constraints
%        during training (src/models/joint_optimizer.py:287-290).

\begin{equation}
\label{eq:stability_margin}
\rho(\tau) = |\lambda_{\min}(0)| - \sum_{i=1}^{n_g} K_i \cdot \frac{\tau_i}{\tau_{\max}}
\end{equation}

\noindent where:
\begin{itemize}
    \item $\rho(\tau)$: Stability margin as a function of delay (positive $\Rightarrow$ stable)
    \item $\lambda_{\min}(0)$: Minimum eigenvalue of the system Jacobian at zero delay
    \item $K_i$: Learnable coupling constant for generator $i$ (sensitivity to delay)
    \item $\tau_i$: Communication delay for generator $i$ (milliseconds)
    \item $\tau_{\max}$: Maximum tolerable delay (typically 500ms for grid control)
    \item $n_g$: Number of generators in the power system
\end{itemize}

\textbf{Physical interpretation:} The stability margin decreases linearly with normalized delay. When $\rho(\tau) > 0$, the closed-loop system remains stable; instability occurs when $\rho(\tau) \leq 0$.


% ============================================================
% SECTION: DUAL-DOMAIN GNN ENCODERS
% ============================================================

% ------------------------------------------------------------
% ENERGY DOMAIN GNN
% ------------------------------------------------------------
% WHAT: Graph neural network encoding power system bus features.
% WHY:  Captures electrical topology and power flow patterns that
%       determine stability characteristics.
% WHERE: First stage of JointOptimizer (src/models/gnn.py:156-254).
% HOW:  Multi-layer GAT with residual connections and layer normalization.

\begin{equation}
\label{eq:energy_gnn}
h_E^{(\ell)} = \text{LayerNorm}\left( h_E^{(\ell-1)} + \text{GAT}^{(\ell)}\left(h_E^{(\ell-1)}, \mathcal{E}_E\right) \right)
\end{equation}

\noindent where:
\begin{itemize}
    \item $h_E^{(\ell)} \in \mathbb{R}^{N \times d}$: Energy node embeddings at layer $\ell$
    \item $h_E^{(0)} = W_{\text{proj}} x_E$: Initial projection from input features $x_E \in \mathbb{R}^{N \times 5}$
    \item $x_E = [P, Q, V, \theta, \omega]$: Per-bus features (active/reactive power, voltage magnitude/angle, frequency)
    \item $\mathcal{E}_E$: Edge set of power grid topology (transmission lines)
    \item $\text{GAT}$: Graph Attention Network with multi-head attention
    \item $N$: Number of buses; $d$: Embedding dimension (default 128)
\end{itemize}

% ------------------------------------------------------------
% COMMUNICATION DOMAIN GNN
% ------------------------------------------------------------
% WHAT: Graph neural network encoding communication network features.
% WHY:  Models latency, bandwidth, and routing that affect control signal delivery.
% WHERE: Parallel to EnergyGNN (src/models/gnn.py:257-365).

\begin{equation}
\label{eq:comm_gnn}
h_I^{(\ell)} = \text{LayerNorm}\left( h_I^{(\ell-1)} + \text{GAT}^{(\ell)}\left(h_I^{(\ell-1)}, \mathcal{E}_I\right) \right)
\end{equation}

\noindent where:
\begin{itemize}
    \item $h_I^{(\ell)} \in \mathbb{R}^{N \times d}$: Communication node embeddings at layer $\ell$
    \item $h_I^{(0)} = W_{\text{proj}} x_I$: Initial projection from input features $x_I \in \mathbb{R}^{N \times 3}$
    \item $x_I = [\tau, R, B]$: Per-node features (delay, data rate, bandwidth)
    \item $\mathcal{E}_I$: Edge set of communication network topology
\end{itemize}

\textbf{Architectural note:} Both GNNs use identical architecture (3 layers, 4 attention heads) but operate on different graph topologies, enabling domain-specific feature extraction.


% ============================================================
% SECTION: PHYSICS-CONSTRAINED ATTENTION
% ============================================================

% ------------------------------------------------------------
% PHYSICS MASK (Impedance-Weighted)
% ------------------------------------------------------------
% WHAT: Attention bias based on electrical distance between buses.
% WHY:  Electrically close buses have stronger interactions; attention
%       should reflect physical coupling strength.
% WHERE: Applied to cross-domain attention (src/models/attention.py:40-60).
% HOW:  Negative impedance ratio creates soft bias (not hard mask).

\begin{equation}
\label{eq:physics_mask}
M_{\text{phys}}[i,j] = -\gamma \cdot \frac{Z_{ij}}{Z_{\max}}
\end{equation}

\noindent where:
\begin{itemize}
    \item $M_{\text{phys}} \in \mathbb{R}^{N \times N}$: Physics-informed attention mask
    \item $Z_{ij}$: Electrical impedance between bus $i$ and bus $j$ (ohms)
    \item $Z_{\max} = \max_{i,j} Z_{ij}$: Maximum impedance for normalization
    \item $\gamma$: Mask strength hyperparameter (default 1.0)
\end{itemize}

\textbf{Effect:} Low impedance (close buses) $\Rightarrow$ $M_{\text{phys}} \approx 0$ (attend strongly). High impedance (distant buses) $\Rightarrow$ $M_{\text{phys}} < 0$ (attend weakly after softmax).

% ------------------------------------------------------------
% CAUSAL MASK (DAG-Based)
% ------------------------------------------------------------
% WHAT: Attention mask enforcing causal ordering in control DAG.
% WHY:  Information should flow from causes to effects, not vice versa.
%       Prevents spurious correlations from reverse causation.
% WHERE: Applied to self-attention on energy domain (src/models/attention.py:106-126).
% HOW:  Uses transitive closure of DAG to compute ancestor relationships.

\begin{equation}
\label{eq:causal_mask}
M_{\text{causal}}[i,j] = \begin{cases}
    0 & \text{if } j \in \text{Ancestors}(i) \text{ or } j = i \\
    -\infty & \text{otherwise}
\end{cases}
\end{equation}

\noindent where:
\begin{itemize}
    \item $M_{\text{causal}} \in \mathbb{R}^{N \times N}$: Causal attention mask
    \item $\text{Ancestors}(i)$: Set of all nodes that causally precede node $i$ in the control DAG
    \item The $-\infty$ values become 0 after softmax, blocking non-causal attention
\end{itemize}

\textbf{Construction:} Computed via transitive closure of the DAG adjacency matrix in $O(N^3)$ time, cached for efficiency.


% ============================================================
% SECTION: CROSS-DOMAIN ATTENTION MECHANISM
% ============================================================

% ------------------------------------------------------------
% SCALED DOT-PRODUCT ATTENTION WITH MASK
% ------------------------------------------------------------
% WHAT: Standard transformer attention augmented with domain-specific masks.
% WHY:  Enables selective information fusion respecting physical constraints.
% WHERE: Core attention computation (src/models/attention.py:230-259).

\begin{equation}
\label{eq:attention}
\text{Attn}(Q, K, V; M) = \text{softmax}\left( \frac{QK^\top}{\sqrt{d_k}} + M \right) V
\end{equation}

\noindent where:
\begin{itemize}
    \item $Q \in \mathbb{R}^{B \times H \times N \times d_k}$: Query tensor (from energy domain)
    \item $K, V \in \mathbb{R}^{B \times H \times N \times d_k}$: Key and Value tensors (from communication domain)
    \item $M \in \mathbb{R}^{N \times N}$: Attention mask ($M_{\text{phys}}$ or $M_{\text{causal}}$)
    \item $d_k = d / H$: Per-head dimension
    \item $B$: Batch size; $H$: Number of attention heads (default 8)
\end{itemize}

% ------------------------------------------------------------
% CROSS-DOMAIN ATTENTION (Energy → Communication)
% ------------------------------------------------------------
% WHAT: Energy queries attending to communication keys/values.
% WHY:  Allows energy domain to selectively gather relevant communication
%       information while respecting electrical proximity.
% WHERE: CrossDomainAttention module (src/models/attention.py:311-366).

\begin{equation}
\label{eq:cross_attention}
h_{\text{cross}} = \text{Attn}(W_Q h_E, W_K h_I, W_V h_I; M_{\text{phys}})
\end{equation}

\noindent where:
\begin{itemize}
    \item $h_{\text{cross}} \in \mathbb{R}^{B \times N \times d}$: Cross-attended embeddings
    \item $W_Q, W_K, W_V \in \mathbb{R}^{d \times d}$: Learnable projection matrices
    \item $h_E$: Energy embeddings (queries)
    \item $h_I$: Communication embeddings (keys and values)
\end{itemize}


% ============================================================
% SECTION: HIERARCHICAL ATTENTION FUSION
% ============================================================

% ------------------------------------------------------------
% TWO-STAGE ATTENTION WITH FUSION
% ------------------------------------------------------------
% WHAT: Complete attention module combining self-attention and cross-attention.
% WHY:  Self-attention captures intra-domain patterns; cross-attention
%       captures inter-domain dependencies; fusion combines both.
% WHERE: HierarchicalAttention module (src/models/attention.py:368-448).
% HOW:  Sequential application of causal self-attention, then cross-attention,
%       then concatenation and feed-forward fusion.

\begin{equation}
\label{eq:causal_self_attn}
h_E' = \text{LayerNorm}\left( h_E + \text{Attn}(h_E, h_E, h_E; M_{\text{causal}}) \right)
\end{equation}

\begin{equation}
\label{eq:hierarchical_fusion}
h_{\text{fused}} = \text{LayerNorm}\left( h_{\text{concat}} + \text{FFN}(h_{\text{concat}}) \right)
\end{equation}

\noindent where $h_{\text{concat}} = [h_E' \| h_{\text{cross}}]$ (concatenation along feature dimension), and:
\begin{itemize}
    \item $h_E' \in \mathbb{R}^{B \times N \times d}$: Causally-refined energy embeddings
    \item $h_{\text{fused}} \in \mathbb{R}^{B \times N \times d}$: Final fused representation
    \item $\text{FFN}$: Two-layer feed-forward network with GELU activation
\end{itemize}

\textbf{Key insight:} Masks are applied \textit{separately} (causal to self-attention, physics to cross-attention), not as element-wise product.


% ============================================================
% SECTION: LEARNABLE COUPLING CONSTANTS
% ============================================================

% ------------------------------------------------------------
% LOG-PARAMETERIZED COUPLING
% ------------------------------------------------------------
% WHAT: Per-generator coupling constants learned from data.
% WHY:  Traditional fixed K is overly conservative; learning K from data
%       tightens stability bounds by ~18%.
% WHERE: LearnableCouplingConstants module (src/models/coupling.py:12-56).
% HOW:  Exponential parameterization guarantees positivity: K = exp(κ).

\begin{equation}
\label{eq:coupling_param}
K_i = \exp(\kappa_i), \quad \kappa_i \in \mathbb{R}
\end{equation}

\noindent where:
\begin{itemize}
    \item $K_i > 0$: Coupling constant for generator $i$ (guaranteed positive)
    \item $\kappa_i$: Learnable parameter (unconstrained real number)
    \item Initialization: $\kappa_i = \log(K_{\text{init}})$, where $K_{\text{init}}$ is auto-scaled per grid (see below)
\end{itemize}

\textbf{Auto-scaled initialization:} $K_{\text{init}}$ is computed per grid to ensure a positive stability margin at the maximum delay:
\begin{equation}
\label{eq:k_init_auto}
K_{\text{init}} = \frac{s \cdot |\lambda_{\min}(0)|}{n_g}, \quad s = 0.5 \text{ (safety factor)}
\end{equation}
This replaces the fixed $K=0.1$ used in the conference version, ensuring well-scaled initialization for grids of any size.

\textbf{Training dynamics:} $K$ is further refined during training via gradient descent, typically decreasing from the initial value as the model learns to tighten stability bounds.


% ============================================================
% SECTION: LOSS FUNCTIONS
% ============================================================

% ------------------------------------------------------------
% STABILITY LOSS (Hinge)
% ------------------------------------------------------------
% WHAT: Penalizes negative stability margins.
% WHY:  Ensures model outputs maintain system stability (ρ > 0).
% WHEN: Applied during training to enforce stability constraint.

\begin{equation}
\label{eq:stability_loss}
\mathcal{L}_{\text{stab}} = \mathbb{E}\left[ \max(0, -\rho(\tau)) \right]
\end{equation}

\noindent This hinge loss is zero when stable ($\rho > 0$) and increases linearly with instability margin.

% ------------------------------------------------------------
% COUPLING LOSS (MSE)
% ------------------------------------------------------------
% WHAT: Aligns empirical and theoretical stability margins.
% WHY:  Trains the model to produce margins consistent with Theorem 1.
% WHEN: Applied with weight α to balance against other objectives.

\begin{equation}
\label{eq:coupling_loss}
\mathcal{L}_{\text{couple}} = \left\| \rho_{\text{emp}} - \rho_{\text{theo}} \right\|_2^2
\end{equation}

\noindent where $\rho_{\text{emp}}$ is computed from model predictions and $\rho_{\text{theo}}$ from \eqref{eq:stability_margin}.

% ------------------------------------------------------------
% TOTAL LOSS
% ------------------------------------------------------------
% WHAT: Combined multi-objective loss function.
% WHY:  Balances optimal power flow, communication QoS, stability, and coupling.

\begin{equation}
\label{eq:total_loss}
\mathcal{L} = \mathcal{L}_{\text{OPF}} + \mathcal{L}_{\text{QoS}} + \mathcal{L}_{\text{stab}} + \alpha \mathcal{L}_{\text{couple}}
\end{equation}

\noindent where:
\begin{itemize}
    \item $\mathcal{L}_{\text{OPF}}$: Optimal power flow loss (generation cost minimization)
    \item $\mathcal{L}_{\text{QoS}}$: Quality of service loss (communication latency/throughput)
    \item $\mathcal{L}_{\text{stab}}$: Stability loss from \eqref{eq:stability_loss}
    \item $\mathcal{L}_{\text{couple}}$: Coupling loss from \eqref{eq:coupling_loss}
    \item $\alpha = 1.0$: Coupling loss weight (optimal from ablation study)
\end{itemize}


% ============================================================
% SECTION: INFORMATION-THEORETIC SEPARATION (OBSERVATION 1)
% ============================================================

% ------------------------------------------------------------
% MUTUAL INFORMATION BOUND
% ------------------------------------------------------------
% WHAT: Empirical observation on information separation between domains.
% WHY:  Ensures domain-specific embeddings remain separable while
%       coupling through learned K, not through spurious correlations.
% WHERE: Validated empirically (I ≈ 10⁻⁴ nats in experiments).
% NOTE: Renamed from "Theorem 2" (conference) to "Observation 1" (journal).

\begin{equation}
\label{eq:mutual_info}
I(h_E; h_I) \leq \epsilon, \quad \epsilon \ll 1
\end{equation}

\noindent where:
\begin{itemize}
    \item $I(h_E; h_I)$: Mutual information between energy and communication embeddings
    \item $\epsilon \approx 10^{-4}$ nats: Empirically observed upper bound
\end{itemize}

\textbf{Interpretation:} Low MI indicates the two domains encode complementary (not redundant) information, validating the dual-encoder architecture.


% ============================================================
% ALGORITHM 1: JointOptimizer Training
% ============================================================
% Lines: 16 (within 22-line limit)
% Source: run_baseline_comparison.py, src/models/joint_optimizer.py

\begin{algorithm}[t]
\caption{JointOptimizer: Energy-Information Co-Optimization}
\label{alg:joint_optimizer}
\begin{algorithmic}[1]
\REQUIRE Power grid graph $\mathcal{G}_E = (V, \mathcal{E}_E)$, Communication graph $\mathcal{G}_I = (V, \mathcal{E}_I)$, delays $\tau$, epochs $T$
\ENSURE Trained parameters $\theta$, learned coupling constants $K$
\STATE Initialize $\theta$ randomly; $K_{\text{init}} \leftarrow s \cdot |\lambda_{\min}|/n_g$; $\kappa \leftarrow \log(K_{\text{init}}) \cdot \mathbf{1}_{n_g}$ \COMMENT{Auto-scaled}
\FOR{epoch $= 1$ to $T$}
    \FOR{each batch $(x_E, x_I, \tau, \lambda_{\min})$}
        \STATE $h_E \leftarrow \text{EnergyGNN}(x_E, \mathcal{G}_E)$ \COMMENT{Encode power system state}
        \STATE $h_I \leftarrow \text{CommGNN}(x_I, \mathcal{G}_I)$ \COMMENT{Encode communication state}
        \STATE $h_E' \leftarrow h_E + \text{CausalAttn}(h_E; M_{\text{causal}})$ \COMMENT{Causal self-attention}
        \STATE $h_{\text{cross}} \leftarrow \text{CrossAttn}(h_E', h_I; M_{\text{phys}})$ \COMMENT{Physics-masked cross-attention}
        \STATE $h_{\text{fused}} \leftarrow \text{FFN}([h_E' \| h_{\text{cross}}])$ \COMMENT{Fusion}
        \STATE $(P^*, Q^*) \leftarrow \text{Decoder}(h_{\text{fused}})$ \COMMENT{Predict optimal dispatch}
        \STATE $K \leftarrow \exp(\kappa)$; \quad $\rho \leftarrow |\lambda_{\min}| - \sum_i K_i \tau_i / \tau_{\max}$ \COMMENT{Theorem 1}
        \STATE $\mathcal{L} \leftarrow \mathcal{L}_{\text{OPF}} + \mathcal{L}_{\text{QoS}} + \mathcal{L}_{\text{stab}} + \alpha \|\rho_{\text{emp}} - \rho\|^2$
        \STATE $\theta, \kappa \leftarrow \text{Adam}(\nabla_{\theta,\kappa} \mathcal{L})$ \COMMENT{Update all parameters}
    \ENDFOR
\ENDFOR
\RETURN $\theta$, $K = \exp(\kappa)$
\end{algorithmic}
\end{algorithm}

\textbf{Complexity:} $O(N^2 \cdot d)$ per batch due to attention; $O(|\mathcal{E}| \cdot d)$ for GNN message passing. Validated on IEEE 14, 30, 39, 57, and 118-bus systems (GPU training).
